% !TeX root = ../se200_referential_regimes.tex

% =============================================================================
\section{Transformation-Invariant Identity}
\label{sec:algebra}
% =============================================================================

The previous section fixed the referents the substrate must preserve.
This section analyzes the conditions under which two representations
of such a referent count as the same entity.
The analysis treats identity as \emph{transformation-invariance}:
an identity regime is determined by the behavior of a declared identity basis
under the transformation basis fixed in Section~\ref{sec:formal_setting}.

The governing idea is familiar from formal ontology:
categories are individuated by identity conditions, not by classificatory convenience,
and properties that a thing can gain or lose cannot carry identity~\citep{guarino1998roles,guarinowelty2002}.
This paper applies that idea at the substrate level and in operational form.
A referent remains the same under a regime when the transformations applied
to its representations do not break the identity basis for that regime.
Two representations may be the same under one regime and different under another.
A substrate that leaves the regime implicit has not fixed reference.
It has deferred the identity question to a later interpretive choice.
Under persistent disagreement, that choice cannot be assumed shared.

% -----------------------------------------------------------------------------
\subsection{Identity Bases}
\label{subsec:identity_bases}
% -----------------------------------------------------------------------------

\begin{definition}[Identity basis]
\label{se200.def.IdentityBasis}
An \emph{identity basis} is the feature, structure, role in reference, or persistence condition
whose preservation constitutes remaining the same referent under a given identity regime.
\end{definition}

An identity basis is the condition a regime tracks through transformation.
The lower-bound construction uses responsibility-bearing identity
for obligation-bearers,
temporal-realization and provenance identity for occurrences,
and descriptive-record identity for records.
It also uses three witnessed split cases.
A rule may be content-fixed or structure-fixed.
A scope may be extension-fixed or structure-fixed.
A plain referent may be locus-fixed or object-fixed.

The inventory supplies carriers.
The identity basis determines sameness for a carrier.
The transformation classification determines whether that basis is preserved or broken.
The lower-bound argument begins when two bases for the same carrier kind
classify an ordinary transformation differently.

% -----------------------------------------------------------------------------
\subsection{Profiles and Collapse}
\label{subsec:profiles_and_collapse}
% -----------------------------------------------------------------------------

The formal setting already defined identity regimes as classification profiles over the transformation basis.
We now use the induced identity relation $\sim_\tau$ to compare regimes.

\begin{definition}[Collapse]
\label{se200.def.Collapse}
Two identity regimes $\tau$ and $\tau'$ \emph{collapse} when they induce the same identity relation:
\[
\sim_\tau \ =\ \sim_{\tau'} .
\]
They are \emph{distinct} when some pair of representations is identified under one regime and separated under the other.
\end{definition}

\begin{definition}[Underdetermined kind]
\label{se200.def.underdetermined_kind}
A reference kind is \emph{underdetermined} when its kind label
does not determine a single identity basis and admissible identity bases
for that kind classify some transformation family differently.
Equivalently, one admissible basis for the kind
classifies a transformation as non-breaking,
while another admissible basis for the same kind
classifies that transformation as breaking.
\end{definition}

Underdetermination is the engine of the result.
When basis plurality is in play for an underdetermined kind,
no single regime profile serves all required bases.
The kind must be represented through distinct identity regimes for those bases.
The split is forced by transformation behavior, not by the name of the kind.

% -----------------------------------------------------------------------------
\subsection{Carrier Kinds Not Split in the Core Construction}
\label{subsec:core_unsplit_kinds}
% -----------------------------------------------------------------------------

Three carrier kinds contribute one regime each to the core.
They are not split by the witness operations used in this construction:
obligation-bearers, occurrences, and records.

An \emph{obligation-bearer} is individuated by persistence as the party that can bear
responsibility, duty, liability, authority, or accountability.
Re-expression, annotation, re-contextualization, reassignment, branching,
provenance extension, and state evolution do not alter that identity basis.
Refinement and decomposition preserve identity when they add detail or decompose descriptions
of the bearer without replacing the bearer.
Substitution breaks identity because it replaces the bearer.

An \emph{occurrence} is individuated by temporal realization and provenance.
Re-expression, annotation, refinement, re-contextualization, and provenance extension do not replace the occurrence.
State evolution preserves identity when it tracks phases or stages of the same occurrence.
Decomposition breaks identity when it produces sub-occurrences with their own temporal individuation.
Substitution breaks identity because it replaces the occurrence.
Branching and reassignment are inapplicable to a fixed realization event.

An \emph{record} is individuated by descriptive persistence without substrate-level causal or normative commitment.
Annotation, refinement, decomposition, branching, and provenance extension
preserve record identity because the record survives enrichment, derivation, copying, and descriptive elaboration.
Re-expression, re-contextualization, reassignment, and state evolution are
identity-neutral for the record as such.
Substitution breaks identity because it replaces the record.

These three carrier-basis pairs are included once in the core.
We name their regimes:
\[
\OBL,\qquad \OCC,\qquad \REC .
\]

% -----------------------------------------------------------------------------
\subsection{Witnessed Split Kinds}
\label{subsec:witnessed_split_kinds}
% -----------------------------------------------------------------------------

Three reference kinds have witnessed split pairs in the lower-bound construction.
Each admits at least the two admissible readings used here,
and an ordinary transformation separates those readings.
The pattern is the same in all three cases.
A kind admits two bases $B_1$ and $B_2$.
A transformation $f$ is non-breaking for $B_1$ and breaking for $B_2$.
When admissible frameworks require both bases,
the kind cannot be governed by a single identity regime.

\begin{lemma}[Split pattern]
\label{se200.lem.SplitPattern}
Let $K$ be a reference kind with two admissible identity bases $B_1$ and $B_2$.
Suppose there is a transformation family $f\in\Fset$ and representations $r,s$
such that $f$ preserves identity under $B_1$ but breaks identity under $B_2$.
If a neutral substrate must support both bases for $K$,
then $K$ requires at least two distinct identity regimes.
\end{lemma}

\begin{proof}
If one regime governed both bases, it would assign a single classification to $f$
for the kind.
But the $B_1$ reading requires $f$ to be non-breaking,
while the $B_2$ reading requires $f$ to be breaking.
No single classification can satisfy both readings.
Thus the kind is underdetermined at $f$.

Applying $f$ to the witnessing representations identifies them under $B_1$
and separates them under $B_2$.
The two bases therefore induce distinct identity relations.
When both bases must be supported, the substrate requires at least two
distinct identity regimes for $K$.
\end{proof}

% -----------------------------------------------------------------------------
\subsection{Scope Under Decomposition}
\label{subsec:scope_split}
% -----------------------------------------------------------------------------

A scope may be individuated by its extension or by its structure.
An extension-fixed scope is individuated by the
cases, persons, places, entities, or conditions it covers.
A structure-fixed scope is individuated
by the internal organization through which that coverage is represented.

\begin{proposition}[Scope Underdetermination]
\label{se200.prop.ScopeUnderdetermination}
The scope reference kind is underdetermined at aggregation/decomposition $\mathrm{AD}$.
\end{proposition}

\begin{proof}
Let $r$ be a scope fixed by extension.
Let $s$ be a decomposition of $r$ into subscopes whose union covers the same cases.
The extension is preserved, so $\mathrm{AD}$ is non-breaking for the extension-fixed reading.
Let $r'$ be a scope fixed by structure.
Let $s'$ be a decomposition of $r'$ into a finer internal partition.
The structure is changed, so $\mathrm{AD}$ is breaking for the structure-fixed reading.
Both readings are admissible scopes.
No single classification of $\mathrm{AD}$ can serve both.
By Lemma~\ref{se200.lem.SplitPattern}, when both bases must be supported,
scope requires at least two identity regimes.
\end{proof}

This yields:
\[
\SCOPEE,\qquad \SCOPES .
\]
$\SCOPEE$ is the extension-fixed scope regime.
$\SCOPES$ is the structure-fixed scope regime.
After decomposition, two representations
can apply to the same cases while no longer being the same scope structure.
State evolution is inapplicable to both scope regimes:
a scope is an applicability domain,
not an enduring referent that passes through internal states.

% -----------------------------------------------------------------------------
\subsection{Rule Under Refinement}
\label{subsec:rule_split}
% -----------------------------------------------------------------------------

A rule may be individuated by its content or by its structure.
A content-fixed rule is individuated by its requirements,
permissions, prohibitions, authorizations, or definitions.
A structure-fixed rule is individuated by the organization of the
text, sections, clauses, delegations, or internal rule structure
through which that content is expressed.

\begin{proposition}[Rule Underdetermination]
\label{se200.prop.RuleUnderdetermination}
The rule reference kind is underdetermined at refinement $\mathrm{RF}$.
\end{proposition}

\begin{proof}
Let $r$ be a rule fixed by normative content.
Let $s$ be a refinement of $r$ that preserves the rule's requirements
while adding detail or reorganizing expression.
The content is preserved, so $\mathrm{RF}$ is non-breaking for the content-fixed reading.
Let $r'$ be a rule fixed by textual or internal structure.
Let $s'$ be a refinement of $r'$ that changes that structure.
The structure is changed, so $\mathrm{RF}$ is breaking for the structure-fixed reading.
Both readings are admissible rules.
No single classification of $\mathrm{RF}$ can serve both.
By Lemma~\ref{se200.lem.SplitPattern}, when both bases must be supported,
rule requires at least two identity regimes.
\end{proof}

This yields:
\[
\RULEC,\qquad \RULES .
\]
$\RULEC$ is the content-fixed rule regime.
$\RULES$ is the structure-fixed rule regime.
After refinement, two representations
can require the same thing while no longer being the same text or rule structure.
State evolution is inapplicable to both rule regimes:
a rule is a normative structure, not an enduring referent that passes through internal states.
Changes to rule content are amendments or refinements,
not state transitions of a persisting object.

% -----------------------------------------------------------------------------
\subsection{Plain Referent Under Forking}
\label{subsec:plain_referent_split}
% -----------------------------------------------------------------------------

A plain referent may be individuated by locus or by object.
A locus-fixed referent is individuated by the
place, site, slot, location, or institutional position that remains fixed.
An object-fixed referent is individuated by the
object, asset, instrument, sample, device, parcel, or thing occupying that locus.

\begin{proposition}[Plain-referent Underdetermination]
\label{se200.prop.PlainReferentUnderdetermination}
The plain-referent kind is underdetermined at branch/fork $\mathrm{BF}$.
\end{proposition}

\begin{proof}
Let $r$ be a plain referent fixed by locus.
Let $s$ be a forked continuation over the same locus.
The locus is preserved, so $\mathrm{BF}$ is non-breaking for the locus-fixed reading.
Let $r'$ be a plain referent fixed by object.
Let $s'$ be a fork into non-identical object continuations.
The object identity is not preserved across both continuations,
so $\mathrm{BF}$ is breaking for the object-fixed reading.
Both readings are admissible plain referents.
No single classification of $\mathrm{BF}$ can serve both.
By Lemma~\ref{se200.lem.SplitPattern}, when both bases must be supported,
plain referents require at least two identity regimes.
\end{proof}

This yields:
\[
\LOC,\qquad \OBJ .
\]
$\LOC$ is the locus-fixed plain-referent regime.
$\OBJ$ is the object-fixed plain-referent regime.
After a fork, two representations
can be the same place while no longer being the same object.
State evolution applies to both plain-referent regimes and preserves identity:
a locus or an object is an enduring referent that persists through changes in its state.

% -----------------------------------------------------------------------------
\subsection{Result of the Algebra}
\label{subsec:algebra_result}
% -----------------------------------------------------------------------------

The reference inventory supplied six carrier kinds.
Three carrier-basis pairs contribute one regime each:
\[
\OBL,\qquad \OCC,\qquad \REC .
\]
Three carrier kinds have witnessed split pairs:
\[
\LOC,\ \OBJ,\qquad
\SCOPEE,\ \SCOPES,\qquad
\RULEC,\ \RULES .
\]
Thus the transformation analysis yields the nine-regime lower-bound core:
\[
\begin{aligned}
&\OBL,\quad \OCC,\quad \REC,\quad \LOC,\quad \OBJ,\\
&\SCOPEE,\quad \SCOPES,\quad \RULEC,\quad \RULES .
\end{aligned}
\]

The additional regimes arise within carrier kinds already named.
They are identity bases separated by ordinary transformations.
This is why counting carrier kinds undercounts the structure needed
for stable reference.
The next section collects the nine core regimes in a single classification table.
